\documentstyle[%


righttag%


,11pt%


,amssymb%


,russian%               remove in Eng.


,izvru%                 Set izven% in Eng.


% ,izvru-ks             for brief communications


]{amsart}





% 





\newcommand{\Ch}{\operatorname{Ch}}


\newcommand{\tts}[1]{}





\theoremstyle{plain}


\theorembodyfont{\it}


\newtheorem{thmnonum}{\hskip\parindent Теорема}


\renewcommand{\thethmnonum}{}





\theoremstyle{definition}


\theorembodyfont{\rm}


\newtheorem{cornonum}{\hskip\parindent Следствие}


\newtheorem{defnnonum}{\hskip\parindent Определение}


\renewcommand{\thecornonum}{}


\renewcommand{\thedefnnonum}{}





%\hoffset=-15mm%                  For Eng.


\setcounter{page}{34}


\god{2000}


\nomer{9~(460)} %                        Remove in Eng.


%\nomer{Vol.\,44, No.\,9}%               Set in Eng.


%\str{\pageref{begin}--\pageref{end}}%  Set in Eng.


\udk{515.124.4}





\author{\it Е.Н.\,Сосов}


% NB:   ^^ remove in Eng.


\title{Об $\size{14}{18pt}\selectfont\varepsilon$-ядре ограниченного


множества в специальном метрическом пространстве }


\thanks{Работа выполнена при финансовой поддержке Российского фонда


фундаментальных исследований, грант \No\,00-01-00308.}





\date{}


%\pagestyle{myheadings}%         Set in Eng.


\pagestyle{plain} %              Remove in Eng.


\begin{document}


%\thispagestyle{plain}%          Set in Eng.


\maketitle


\label{begin}








В данной статье доказывается, что отображение, ставящее в соответствие каждому 


непустому ограниченному множеству в специальном метрическом пространстве его    


$\varepsilon$-ядро, является равномерно непрерывным относительно метрики


Хаусдорфа. Это свойство обобщает свойство сильной устойчивости чебышевского


центра ограниченного множества в равномерно выпуклом банаховом пространстве [1].





\section{Необходимые определения и теоремы}





Напомним некоторые определения и понятия для 


метрического пространства $(X,\rho)$. Чебышевским 


радиусом ограниченного множества $M \subset X$ называется число [2]


$$


R(M) = \inf_{у \in X}\sup_{x\in M}{\rho (x,y)}.


$$


Точка $z \in X$, для которой


$$


\sup_{x \in M}{\rho (x,z)}=\inf_{у \in X}{\sup\limits_{x \in M}{\rho (x,y)}},


$$


называется чебышевским центром множества $M$ [2]. 


Метрикой Хаусдорфа на множестве 


$B[X]$ всех непустых ограниченных замкнутых множеств 


пространства $X$ называется метрика ([3], с.\,223) 


$$


\delta (M,N) =


\max\{\sup_{x \in M}{\rho(x,N)};\ \sup\limits_{y \in


N}{\rho(y,M)}\}


$$


для $M, N \in B[X].$


Очевидно, что на множестве $B(X)$ всех непустых ограниченных множеств


пространства $X$ это выражение определяет псевдометрику Хаусдорфа.





Пусть $\varepsilon > 0$, $M \in B(X)$,


$B[x,r]$ $(B(x,r))$ --- замкнутый (открытый)


шар с центром в $x \in X$


радиуса $r > 0$ и $B(M,r) = \bigcup \{B(x,r) : x \in M\}$.


Введем множества





$K_{\varepsilon}(M) = \{y


\in X : \sup\limits_{x \in M}{\rho(x,y)} \leq R(M) + \varepsilon \},$





$\varepsilon (M) = \bigcap \{B[x,R(M)


+ \varepsilon ] : M \subset B[x,R(M) + \varepsilon ]\},$





$k_{\varepsilon}(M) = \bigcap \{B[x,r(x,M) +


\varepsilon ] : x \in \varepsilon(M)\},$


где $r(x,M) = (\rho(x,M) + \sup\limits_{y \in M}{\rho(x,y)})/2.$





\begin{defnnonum}


Назовем $\varepsilon$-ядром множества $M$ множество $K_{\varepsilon }(M)$,


где $M \in B(X)$ и $0 < \varepsilon < R(M)$.


\end{defnnonum}


 


Отметим некоторые элементарные свойства введенных множеств, доказательства


которых нетрудно получить непосредственно из определений.





1. Для каждого множества $M \in B(X)$ множества $K_{\varepsilon}(M)$, 


$\varepsilon (M)$ принадлежат $B[X]$ и могут быть 


представлены в следующих формах:


$\varepsilon (M) = \bigcap \{B[x,R(M) +


\varepsilon ] : x \in K_{\varepsilon}(M)\}$,


$K_{\varepsilon}(M) = \bigcup \{y \in X : M \subset B[y,R(M)


+ \varepsilon] \} = 


\bigcup \{y \in X : \varepsilon (M) \subset B[y,R(M) + \varepsilon] \} =


\bigcap \{B[x,R(M) + \varepsilon ] : x \in \varepsilon(M)\}$.





2. Если $\alpha \leq \varepsilon $,


то $K_{\alpha }(M) \subset K_{\varepsilon}(M) $.





3. Если для множества $M \in B(X)$ существует чебышевский центр, то


$\bigcap \{K_{\varepsilon}(M) : \varepsilon > 0 \}$ ---


множество всех чебышевских центров этого множества.





4. Для каждого множества $M \in B(X)$ имеют место неравенства 


$R(K_{\varepsilon}(M)) \leq R(M) + \varepsilon $, 


$R(M) \leq R(\varepsilon (M)) \leq R(M) + \varepsilon $.





5. Пусть каждый замкнутый шар пространства $X$ является выпуклым 


множеством в следующем смысле: вместе со всякими двумя точками


$x$, $y$ он содержит каждую точку $z$, удовлетворяющую условиям


$\rho (x,z) = \rho (y,z) = \rho (x,y)/2$. Тогда множества


$K_{\varepsilon}(M)$, $\varepsilon (M)$, $k_{\varepsilon}(M)$


являются выпуклыми в том же смысле для каждого $M \in B(X)$.





В дальнейшем на метрическое пространство $X$ будем 


налагать следующее условие [4]:


для любых $x$, $y$ из $X$ найдется точка $z \in X$ такая, что 


\begin{equation}


\rho (x,z) = \rho (y,z) = \rho (x,y)/2.


\end{equation}





Сформулируем полученные в данной статье результаты. 





\begin{lem}


Имеет место неравенство 


$\delta (B[x,r],B[y,R]) \leq \rho (x,y) + |R - r|$ для всех $x$, $y$ из 


$X$ и $R >0$, $r > 0$.


\end{lem}





\begin{lem}


Пусть $M_{\alpha n} \in B(X)$, $N_{\alpha n} \in B(X)$ для каждого 


натурального $n$ и каждого $\alpha $ из некоторого множества


индексов $A$. Если


$\bigcap\limits_{\alpha \in A}M_{\alpha n}\neq \emptyset$,


$\bigcap \limits_{\alpha \in A}N_{\alpha n}\neq \emptyset$


для каждого натурального $n$ и


$\delta (M_{\alpha n},N_{\alpha n}) \to 0$


$(n \to \infty)$ для каждого


$\alpha \in A$, то


$$


\delta \Big(\bigcap\limits_{\alpha \in A}M_{\alpha n},


\bigcap \limits_{\alpha \in A}N_{\alpha n}\Big) \to 0


\ \; (n \to \infty ).


$$


\end{lem}





\begin{thmnonum}


Пусть для каждого натурального $n$ непустые


множества $M_{n}$, $N_{n}$ ограничены и 


$\delta (M_{n},N_{n}) \to 0$ $(n \to \infty)$. Тогда





$\mathrm A.$ $\delta (K_{\varepsilon }(M_{n}),K_{\varepsilon }(N_{n}))


\to 0$ $(n \to \infty)$,





$\mathrm B.$ $\delta (\varepsilon (M_{n}),\varepsilon (N_{n}))


\to 0$ $(n \to \infty )$,





$\mathrm C.$ Если для каждого натурального $n$ множества 


$k_{\varepsilon }(M_{n})$, $k_{\varepsilon }(N_{n})$ непустые, то 


$\delta (k_{\varepsilon }(M_{n}),\linebreak k_{\varepsilon }(N_{n})) \to 0$


$(n \to \infty) $.


\end{thmnonum}





\begin{note}


Утверждения теоремы равносильны тому, что отображения 


$K_{\varepsilon } : B(X) \to B(X)$,


$\varepsilon : B(X) \to B(X)$,


$k_{\varepsilon } : \wt B(X) \to B(X)$, где $ \wt B(X)$ --- подпространство 


таких элементов $M \in B(X)$,


что $k_{\varepsilon }(M) \neq \emptyset $, равномерно


непрерывны относительно псевдометрики Хаусдорфа.


\end{note}





Из свойства 3, теоремы и леммы 2 следует





\begin{cornonum}


Пусть для каждого натурального $n$ множества $M_{n} \in B(X)$,


$N_{n} \in B(X)$ обладают непустыми множествами чебышевских


центров $\Ch(M_{n})$, $\Ch(N_{n})$ и 


$\delta (M_{n},N_{n}) \to 0$ ($n \to \infty $). Тогда 


$\delta (\Ch(M_{n}),\Ch(N_{n})) \to 0$ ($n \to \infty $).


\end{cornonum}





\begin{note}


Это свойство обобщает свойство сильной устойчивости чебышевского


центра ограниченного множества в равномерно выпуклом банаховом


пространстве [1].


\end{note}





\section{Доказательства полученных результатов}





{\bf Доказательство леммы 1}.


Пусть $z \notin B[x,r]$. Тогда из условия (1) следует,


что для каждого $\varepsilon >0$ найдутся $u(\varepsilon ) \in B[x,r]$ и 


$v(\varepsilon ) \notin B[x,r]$, удовлетворяющие следующим условиям:





a) $\rho (z,u(\varepsilon )) + \rho (u(\varepsilon),x) = \rho (z,x)$, 


$\rho (z,B[x,r]) \geq \rho (z,u(\varepsilon )) - \varepsilon $;





b) $\rho (z,v(\varepsilon )) + \rho (v(\varepsilon),x) = \rho (z,x)$, 


$\rho (z,B[x,r]) \leq \rho (z,v(\varepsilon )) + \varepsilon $.





\noindent{}Кроме того, 


$\rho (z,x) - r - \varepsilon \leq \rho (z,x) -


\rho (u(\varepsilon ),x) - \varepsilon = 


\rho (z,u(\varepsilon )) - \varepsilon \leq \rho (z,B[x,r])$, 


$\rho (z,B[x,r]) \leq \rho (z,v(\varepsilon )) + \varepsilon


= \rho (z,x) - \rho (x,v(\varepsilon )) + \varepsilon \leq \rho (z,x) - r 


+ \varepsilon.$


Отсюда в силу произвольности $\varepsilon $ следует, что 


$\rho (z,B[x,r]) = \rho (z,x ) - r$. Используя это равенство, получим


\begin{multline*}


\delta (B[x,r],B[y,R]) = 


\max\{\sup_{z \in B[x,r]}{\rho(z,B[y,R])};


\ \sup_{w \in B[y,R]}{\rho(w,B[x,r])}\} = \\


%


=\max\{\sup_{z \in B[x,r]}{\max\{\rho(z,y) - R,0\}};


\ \sup_{w \in B[y,R]}{\max\{\rho(w,x)- r,0\}}\} \leq \\ 


%


\le \max\{\max\{\rho(x,y) + r - R,0\};\ \max\{\rho(y,x) + R - r,0\}\}


= \rho (x,y) + |R - r|.\quad\square


\end{multline*}





{\bf Доказательство леммы 2}.


Используем сокращенное обозначение $\bigcap = \bigcap


\limits_{\alpha \in A}$. Будем доказывать лемму от противного.


Пусть найдутся подпоследовательности


$\{\bigcap M_{\alpha m}\}$, $\{\bigcap N_{\alpha m}\}$


последовательностей $\{\bigcap M_{\alpha n}\}$, $\{\bigcap 


N_{\alpha n}\}$, натуральное число $m_{0}$ и константа $c > 0$, для которых 


выполняются неравенства


$\delta (\bigcap M_{\alpha m},\bigcap N_{\alpha m}) > c$ при $m > m_{0}$.


Отсюда следует, что найдутся подпоследовательности


$\{\bigcap M_{\alpha l}\}$, $\{\bigcap N_{\alpha l}\}$ 


последовательностей $\{\bigcap M_{\alpha m}\}$,


$\{\bigcap N_{\alpha m}\}$, натуральное число $l_{0}$, для которых


одно из двух выражений под знаком $\max$ в псевдометрике


Хаусдорфа (для определенности, первое выражение) удовлетворяет неравенству


$\sup\{\rho(x_{l},\bigcap N_{\alpha l}):x_{l}\in\bigcap M_{\alpha l}\}>c$


при $l > l_{0}$.


Поэтому найдутся $x_{l} \in \bigcap M_{\alpha l}$ такие, что 


$\rho(x_{l},\bigcap N_{\alpha l})


> c$ при $l > l_{0}$.


Следовательно, $B(x_{l},c/3) \cap B(\bigcap N_{\alpha l},c/3)=\emptyset$


при $l > l_{0}$. Отсюда следует, что найдется такое $\beta \in A$, что 


$B(x_{l},c/3) \cap N_{\beta l} = \emptyset$ при $l > l_{0}$. 


Но тогда


$\delta (M_{\beta l},N_{\beta l}) \geq \sup\{\rho(x_{\beta l},N_{\beta l}): 


x_{\beta l}\in M_{\beta l}\}\ge\rho(x_{\beta l},N_{\beta l}) \geq c/3>0$


при $l > l_{0}$, что противоречит условию леммы 2.$\quad\square$





{\bf Доказательство теоремы 1}. Докажем прежде всего, что


\begin{equation}


|R(M) - R(N)| \leq \delta (M,N)


\end{equation}


для всех $M$, $N$ из $B(X)$ [1].


Пусть $x \in M$ и $\tau  > 0$. Тогда найдется точка $z \in N$ такая, что


$\rho(x,z) \leq \rho(x,N) + \tau $ . Кроме того, получим неравенства 


$\rho(x,y) \leq \rho(x,z) + \rho(z,y) \leq \rho(x,N) +


\rho(z,y) + \tau \leq \sup\limits_{x \in M}{\rho(x,N)}


+ \sup\limits_{z \in N}{\rho(z,y)} + \tau  \leq 


\delta(M,N) + \sup\limits_{z \in N}{\rho(z,y)} + \tau $ для $y \in X$. 


Отсюда следует, что 


$\inf\limits_{y \in X}{\sup\limits_{x \in M}{\rho (x,y)}} \leq


\delta(M,N) + \sup\limits_{z \in N}{\rho(z,y)} + \tau$.


Переходя к точной нижней грани в правой части и используя


произвольность выбора $\tau > 0$, получим


$$


R(M) = \inf\limits_{y \in X}{\sup\limits_{x \in M}{\rho (x,y)}} \leq


\delta(M,N) + \inf\limits_{y \in X}{\sup\limits_{z \in


N}{\rho(z,y)}} = \delta(M,N) + R(N).


$$ 


Неравенство $R(N) \leq \delta(M,N) + R(M)$ получается аналогично.


Пусть утверждение A теоремы неверно. Тогда найдутся 


подпоследовательности $\{M_{m}\}$, $\{N_{m}\}$


последовательностей $\{M_{n}\}$, $\{N_{n}\}$,


натуральное число $m_{0}$ и константа $c > 0$, для которых 


при $m > m_{0}$ выполняется неравенство 


$\delta (K_{\varepsilon }(M_{m}),K_{\varepsilon }(N_{m})) > с$.


Отсюда следует, что найдутся подпоследовательности


$\{M_{l}\}$, $\{N_{l}\}$ последовательностей $\{M_{m}\}$, 


$\{N_{m}\}$, натуральное число $l_{0}$, для которых одно


из двух выражений под знаком $\max$ в псевдометрике


Хаусдорфа (для определенности, первое выражение) удовлетворяет


при $l > l_{0}$ неравенству 


$\sup\{\rho(x_{l},K_{\varepsilon }(N_{l})) : x_{l}


\in K_{\varepsilon }(M_{l})\} > c$. Поэтому при $l > l_{0}$ найдется 


$z_{l} \in K_{\varepsilon }(M_{l})$ такое, что 


$\rho(z_{l},K_{\varepsilon }(N_{l})) > c$. Следовательно, при $l > l_{0}$\;


$B(z_{l},c/3) \cap B(K_{\varepsilon }(N_{l}),c/3) = \emptyset$.


Отсюда и из свойства 1 следует, что при $l > l_{0}$ 


найдется такое $u_{l} \in N_{l}$, что


$B(z_{l},c/3) \cap B[u_{l},R(N_{l}) + \varepsilon ] = \emptyset$. 


По условию теоремы $\delta (M_{l},N_{l}) \to 0$ ($l \to \infty $).


Следовательно, найдутся


элементы $v_{l} \in M_{l}$ такие, что $\rho (u_{l},v_{l}) \to 0$ ($l \to \infty $).


Кроме того,


\begin{multline}


\delta (B[v_{l},R(M_{l}) + \varepsilon ],B[u_{l},R(N_{l}) + \varepsilon]) 


\geq \sup\limits_{w_{l} \in B[v_{l},R(M_{l}) + \varepsilon]}


{\rho(w_{l},B[u_{l},R(N_{l}) + \varepsilon])} \geq \\ 


\ge \rho(z_{l},B[u_{l},R(N_{l}) + \varepsilon]) \geq c/3 >0


\ \; \text{при}\ \; l > l_{0}. 


\end{multline}


Но из леммы 1 и неравенства (2) следует, что первое 


выражение в этих неравенствах стремится к нулю при $l > l_{0}$.


Получили противоречие. Таким образом, утверждение A доказано.





Пусть утверждение B теоремы неверно. Тогда найдутся 


подпоследовательности $\{M_{m}\}$, $\{N_{m}\}$


последовательностей $\{M_{n}\}$, $\{N_{n}\}$,


натуральное число $m_{0}$ и константа $c > 0$, для которых 


при $m > m_{0}$ выполняется неравенство


$\delta (\varepsilon (M_{m}),\varepsilon (N_{m})) > с$.


Отсюда следует, что найдутся подпоследовательности


$\{M_{l}\}$, $\{N_{l}\}$ последовательностей $\{M_{m}\}$, 


$\{N_{m}\}$, натуральное число $l_{0}$, для которых одно


из двух выражений под знаком $\max$ в псевдометрике


Хаусдорфа (для определенности, первое выражение)


удовлетворяет при $l > l_{0}$ неравенству 


$\sup\{\rho(x_{l},\varepsilon (N_{l})) :


x_{l} \in \varepsilon (M_{l})\} > c$. Поэтому при $l > l_{0}$ найдется 


$z_{l} \in \varepsilon (M_{l})$ такое, что


$\rho(z_{l},\varepsilon (N_{l})) > c$. Следовательно, при $l > l_{0}$\;


$B(z_{l},c/3) \cap B(\varepsilon (N_{l}),c/3) = \emptyset$.


Отсюда и из свойства 1 следует, что при $l > l_{0}$ найдется 


такое $u_{l} \in K_{\varepsilon }(N_{l})$, что


$B(z_{l},c/3) \cap B[u_{l},R(N_{l}) + \varepsilon ] = \emptyset$. 


Из утверждения A и условия теоремы следует, что 


$\delta (K_{\varepsilon }(M_{l}),


K_{\varepsilon }(N_{l})) \to 0$($l \to \infty $).


Следовательно, найдутся элементы 


$v_{l} \in K_{\varepsilon }(M_{l})$ такие, что $\rho (u_{l},v_{l}) \to 0$


($l \to \infty $). В этом случае также имеют место неравенства (3). Повторяя 


рассуждение из доказательства утверждения $A$, получаем противоречие.





Пусть утверждение C теоремы неверно. Тогда найдутся 


подпоследовательности $\{M_{m}\}$, $\{N_{m}\}$


последовательностей $\{M_{n}\}$, $\{N_{n}\}$,


натуральное число $m_{0}$ и константа $c > 0$, для которых 


при $m > m_{0}$ выполняется неравенство


$\delta (k_{\varepsilon }(M_{m}),k_{\varepsilon }(N_{m})) > с$.


Отсюда следует, что найдутся подпоследовательности


$\{M_{l}\}$, $\{N_{l}\}$ последовательностей $\{M_{m}\}$, 


$\{N_{m}\}$, натуральное число $l_{0}$, для которых


одно из двух выражений под знаком $\max$ в псевдометрике


Хаусдорфа (для определенности, первое выражение)


удовлетворяет при $l > l_{0}$ неравенству


$\sup\{\rho(x_{l},k_{\varepsilon }(N_{l})) :


x_{l} \in k_{\varepsilon }(M_{l})\} > c$.


Поэтому при $l > l_{0}$ найдется  $z_{l} \in k_{\varepsilon }(M_{l})$


такое, что $\rho(z_{l},k_{\varepsilon }(N_{l})) > c$.


Следовательно, при $l > l_{0}$\;


$B(z_{l},c/3) \cap B(k_{\varepsilon }(N_{l}),c/3) = \emptyset$.


Отсюда и из свойства 1 следует, что при $l > l_{0}$


найдется такое $u_{l} \in \varepsilon (N_{l})$,


что $B(z_{l},c/3) \cap B[u_{l},r(u_{l},N_{l}) + \varepsilon ] = \emptyset$.


Из утверждения B и условия теоремы следует, что 


$\delta (\varepsilon (M_{l}),\varepsilon (N_{l})) \to 0$ 


($l \to \infty $). Следовательно, найдутся элементы 


$v_{l} \in \varepsilon (M_{l})$ такие, что $\rho (u_{l},v_{l}) \to 0$


($l \to \infty $). Кроме того,


\begin{multline}


\delta (B[v_{l},r(v_{l},M_{l}) +


\varepsilon ],B[u_{l},r(u_{l},N_{l}) + \varepsilon ]) 


\geq \sup\limits_{w_{l} \in B[v_{l},r(v_{l},M_{l}) + \varepsilon]}


{\rho(w_{l},B[u_{l},r(u_{l},N_{l}) + \varepsilon])} \geq \\ 


%


\ge \rho(z_{l},B[u_{l},r(u_{l},N_{l})+\varepsilon]) \geq c/3 >0


\ \; \text{при}\ \; l > l_{0}. 


\end{multline}


Но из неравенства


$$|r(u_{l},N_{l}) - r(v_{l},M_{l})| = |(\rho(u_{l},N_{l}) + 


\sup\limits_{w_{l} \in N_{l}}{\rho(u_{l},w_{l})})/2 - (\rho(v_{l},M_{l}) + 


\sup\limits_{w_{l} \in M_{l}}{\rho(v_{l},w_{l})})/2| \leq \rho(u_{l},v_{l})


$$


и леммы 1 следует, что первое выражение в неравенствах (4) стремится к нулю


при $l > l_{0}$. Получили противоречие.$\quad\square$





\begin{thebibliography}{9}





\bibitem{1}


Белобров П.К. {\em О чебышевской точке системы множеств} // Изв.


вузов. Математика. -- 1966. -- \No\,6. -- С.\,18--24.


\bibitem{2}


Гаркави А.Л. {\em О наилучшей сети и наилучшем сечении множеств в


нормированном пространстве} // Изв. АН СССР. Сер.


матем. -- 1962. -- Т.\,26. -- No\,1. -- С.\,87--106.


\bibitem{3}


Куратовский К. {\em Топология}. -- М.: Мир, 1966. -- Т.\,1. -- 594~с.


\bibitem{4}


Ефремович В.А. {\em Неэквиморфность пространств Евклида и Лобачевского} //


УMH. -- 1949. -- Т.4. -- Вып.\,2. -- С.\,178--179.





\end{thebibliography}





\vskip 2cm





\post{\it Казанский государственный}{\it Поступила}


\post{\it педагогический университет}{11.05.1999}





\label{end}


\end{document}


